% Other strands of CEP research reinforce this gap. Bayesian-network extensions \cite{wasserkrug2010efficient, wasserkrug2012model} estimate probabilities for observed or derived events but require the presence of an event instance to initiate propagation. Interval-based approaches to temporal uncertainty \cite{zhang2013recognizing} reinterpret event timestamps but always presume the existence of the event whose timestamp is being constrained. Methods grounded in lineage tracking and probabilistic NFAs \cite{shen2008probabilistic, kawashima2010complex, chuanfei2010complex, wang2013complex} compute confidence values for complex event patterns over stochastic input streams, yet their evaluation mechanisms remain inactive when the stream falls silent. Likewise, frameworks such as CEP2U \cite{cugola2015introducing} and uncertainty-aware extensions of Esper and Flink \cite{moreno2019managing} account for uncertainty in measurements and rule firing while offering no semantics for absent events. Across this body of work, progression only occurs when new inputs are delivered; in the presence of temporary data gaps, processing halts, and no provision is made for preserving the continuity of computation.

% A small number of IoT-oriented studies acknowledge the impact of missing events on event-driven analytics and attempt to address the issue through external data-level reconstruction or synthetic event generation. For example, \citet{gkoulis2025assessing} evaluate ``event fabrication'' methods—such as naive interpolation and exponential smoothing—for filling gaps in sensor data to support real-time IoT applications. Similarly, \citet{gkoulis2024hybrid} incorporate an event-generation mechanism into a quality-aware CEP platform to compensate for missing sensor updates, characterizing the resulting complex events using completeness and timeliness metrics. While these approaches demonstrate the practical need to preserve continuity in IoT event streams, both treat gap-filling as a preprocessing or augmentation step applied \emph{before} or \emph{outside} of the CEP semantics. The event processor remains unaware that an expected observation failed to materialize, cannot distinguish fabricated from genuine inputs, and does not propagate uncertainty associated with the generated values. As a result, these methods do not provide an architectural mechanism for integrating reconstruction with high-level reasoning or maintaining internal-state continuity during data gaps.

% wasserkrug2010efficient - proposed a probabilistic event model using Bayesian networks and a Monte Carlo sampling algorithm to approximate event probabilities in real time.

% wasserkrug2012model - They later extended their previous work by representing both imprecise event data (e.g., noisy readings) and uncertain inference rules, where derived events have associated likelihoods based on evidence and prior probabilities

% zhang2013recognizing - introduced a CEP framework for recognizing event patterns when timestamps are imprecise, modeling each event’s time as an interval rather than a single value. Their approach rewrites pattern queries into constraint-based evaluations that determine whether possible timestamp orderings satisfy the pattern

% shen2008probabilistic – propose a probabilistic stream processor that extends SASE+ (CEP Engine) with lineage based confidence, allowing probability to be calculated over compositive event sequences. Uses Arctive Instance graph (AIG) to track NFA states over streams

% kawashima2010complex - extend the SASE+ CEP engine to support probabilistic event streams by introducing a DFA-based evaluation model that computes confidence values for complex patterns over noisy sensor data.

% moreno2019managing – extends CEP (Ex Esper and Flink) with a Java framework for handling measurement uncertainty, occurrence uncertainty, and uncertainty in rule inference.

% cugola2015introducing - introduced CEP2U, a probabilistic extension to CEP that models both measurement uncertainty and rule uncertainty using probability distributions and Bayesian networks. The framework uses TESLA and T-Rex

% chuanfei2010complex - proposes IPF-DA, a complex event detection algorithm that processes probabilistic event streams in a single scan by storing only relevant event instances in a Chain Instance Queue (CIQ) and pruning those that cannot contribute to a valid complex event.

% wang2013complex - proposes DPCEP, a distributed probabilistic CEP engine that extends SASE with probabilistic NFAs, parallel active-instance stacks, and optimized query-plan evaluation to detect and compute confidences for complex events across multiple distributed event streams